% !TeX root = ../main.tex
\section{Implementation}

\subsection{Prototype and Detector Logic}

\archsyncrepository{archsync-core}{https://github.com/Little-Boy-s-ArchSync/archsync-core} implements the architecture contract, schema and semantic validation, normalized graph comparison, and deny/allow/require/require-path rules. \archsyncrepository{archsync-guardian}{https://github.com/Little-Boy-s-ArchSync/archsync-guardian} implements TypeScript extraction, source evidence, and the Git-diff gate. Guardian delegates conformance semantics to Core so repository and incremental checks use the same rule engine.

For a concrete PostgreSQL example, consider the following code inside the \texttt{order-service} component:
\begin{verbatim}
import { Pool as PgPool } from "pg";
const db = new PgPool({
  connectionString:
    "postgres://postgres:5432/orders"
});
await db.query("select 1");
\end{verbatim}
The import detector records \texttt{PgPool} as a local name for a recognized \texttt{pg} constructor. The initializer pass recognizes \texttt{new PgPool}, resolves the inline \texttt{connectionString}, and associates \texttt{db} with target \texttt{postgres}, derived from the URL hostname. At \texttt{db.query}, the call detector requires that receiver association before emitting the directed edge $(\texttt{order-service},\texttt{data},\texttt{postgres})$. Merely importing \texttt{pg}, constructing an unrelated client, or calling an unrelated object's \texttt{query} method does not satisfy this detector. Construction alone does not emit the database-access edge. Evidence identifies the query call's file, line, column, detector, and excerpt.

The analysis uses per-file identifier maps and four passes over variable declarations; it is not general alias analysis or interprocedural data-flow analysis. Imported client wrappers, arbitrary reassignment, shadowed identifiers, and relationships requiring resolution across files are not comprehensively handled. Endpoint resolution uses literals, mapped variables, and naming conventions such as \texttt{DATABASE\_URL} for \texttt{postgres}. For binary expressions it tries the right-hand expression first; this deterministic heuristic does not evaluate runtime environment values or branch conditions. The PostgreSQL evidence confidence of 0.99 is a fixed detector annotation, not a calibrated probability. Redis operations similarly require a recognized client, and AMQP operations require a channel associated with a recognized connection.

\subsection{Controlled Evaluation Artifact}

\archsyncrepository{archsync-benchmark}{https://github.com/Little-Boy-s-ArchSync/archsync-benchmark} supplies the synthetic Order Platform, 20 executable patches with group-authored ground truth, and the 40-signal positive/hard-negative corpus. Each patch runs on a fresh copy. Git-diff experiments additionally construct a baseline commit, perform cold and warm checks, and compare incremental head observations with a separate full scan. This full scan checks incremental equivalence; it does not independently establish extraction accuracy.

Measurements use \archsynccommit{Core}{https://github.com/Little-Boy-s-ArchSync/archsync-core}{2affbbb0da859a32b9b9079b4bf718fc7b14993b}{2affbbb0da85}, \archsynccommit{Guardian}{https://github.com/Little-Boy-s-ArchSync/archsync-guardian}{8779bf7965b3bd15f25834f17ca5321c85ae3f43}{8779bf7965b3}, and \archsynccommit{Benchmark}{https://github.com/Little-Boy-s-ArchSync/archsync-benchmark}{24d63ebf2fc3075a1d64f1eaff38cdc0b7f586fb}{24d63ebf2fc3}. Later development heads do not replace these measured artifacts. Source/patch manifests, raw results, package hashes, historical verification commands, and the detector/gate audit remain in the supplementary artifact. These controls support inspection, not independent accuracy evidence. \texttt{archsync-mcp} is excluded.
